\chapter{Methods}
\label{ch:methods}

This chapter describes the research design used to answer the research questions
defined in Chapter~\ref{ch:introduction}.
It explains which methods and procedures are used at each stage, why those methods
are appropriate, and what the known limitations and pitfalls of the chosen approach
are.
The thesis addresses a problem that spans software architecture, containerised
deployment, cloud infrastructure, and communication-level security.
Studying these layers in a single undifferentiated experiment would confound the
effects of partition boundary placement, orchestration overhead, and security mechanisms
into one mixed result.
The methodology therefore follows a deliberate \emph{staged experimental pipeline}
in which each stage changes one layer of the system at a time, preserves the model
and benchmark contract, and compares against the nearest relevant baseline.

% ─────────────────────────────────────────────────────────────────────────────
\section{Overall Research Design}
\label{sec:methods-design}

The thesis is organised around two research strands and seven empirical stages plus two
synthesis stages.
RQ1 studies how architectural split choice and deployment context affect
microservice-based neural network inference.
RQ2 studies the performance and operational cost of communication-level and runtime-level
security hardening applied to the AKS deployment path selected from RQ1.

The stages unfold in the following sequence.
RQ1.1 evaluates four coarse ResNet-18 stage boundaries under tightly controlled local
CPU-only execution in order to identify which two-part decompositions are credible
enough to carry forward before the thesis moves to orchestrated or cloud deployment.
RQ1.2 tests one block-level split inside the carried-forward late-stage region to
determine whether finer granularity reveals meaningful differences beyond the coarse
level; it is a targeted refinement, not a new blind search.
RQ1.3 is a hybrid analytical stage that explains the RQ1.1 and RQ1.2 latency patterns
using two structural properties of the partition -- activation-transfer size and compute
distribution -- drawing on the prior benchmark summaries and on a separately collected
internal timing experiment on the monolithic model, without conducting a new split
benchmark.
RQ1.4 moves a predefined family of one-to-five service chains into a local single-node
\texttt{kind} Kubernetes cluster to measure how end-to-end latency accumulates as
synchronous service boundaries are added under in-cluster orchestration.
RQ1.5 re-runs the same predefined chain family on Azure Kubernetes Service (AKS) to
test whether the local Kubernetes ordering and overhead pattern transfer directionally
to managed cloud execution.
RQ2.1 adds managed-Istio mTLS, service accounts, and authorisation policy to the
selected AKS deployment path, measuring the resulting latency, resource, and operational
overhead using a paired execution design.
RQ2.2 adds AMD SEV-SNP VM-level confidential execution to the already
communication-secured chain from RQ2.1, measuring the incremental cost of selective and
full two-service confidential-node placement.
RQ2.3 synthesises the RQ2.1 and RQ2.2 results into a trade-off comparison across cost,
operational complexity, and protection scope, without collecting any new benchmark data.

The common design principle across all stages is to change one layer of the system at a
time.
This principle follows benchmarking guidance that treats isolation of variables as the
primary condition for interpretable systems performance results~\cite{benchmarking_sc15}.

% ─────────────────────────────────────────────────────────────────────────────
\section{System Under Study}
\label{sec:methods-system}

\subsection{Model}

The neural network used in all stages is ResNet-18 with pretrained ImageNet
weights~\cite{resnet}.
ResNet-18 is organised into four residual stages (\texttt{layer1} through
\texttt{layer4}), each of which doubles the number of channels and halves the spatial
resolution of the intermediate feature maps.
This stage structure provides architecturally motivated split candidates: the four stage
boundaries correspond to qualitatively distinct regions of the network with different
activation sizes and different computational loads.

ResNet-18 is chosen for three reasons.
First, its stage structure supports a principled coarse-to-fine screening methodology.
Second, its size enables CPU-only inference at latencies realistic enough for measurement
while avoiding the confounds of GPU scheduling variance in the early controlled stages.
Third, as a standard reference architecture, its behaviour under partition is directly
relatable to prior split-computing and DNN-partitioning work~\cite{neurosurgeon,deeperthings,survey_partitioning}.

All experiments use FP32 precision with a fixed input shape of
$1 \times 3 \times 224 \times 224$ and a batch size of one.
Batch size one reflects the latency-oriented framing of the research questions: the
thesis measures per-request end-to-end inference latency rather than throughput.

\subsection{Microservice Framework}

A proof-of-concept gRPC-based microservice framework was developed for this thesis.
Each service runs as an independent process -- or, in the Kubernetes stages, as an
independent pod -- that receives an input tensor over gRPC, executes its assigned
segment of the ResNet-18 model, and transmits the resulting intermediate activation to
the next service in the chain.
The final service returns the classification output to a benchmark client.

gRPC is chosen as the inter-service protocol because it is a production-grade RPC
framework with binary serialisation via Protocol Buffers, native streaming support, and
broad adoption in cloud-native deployments.
Using gRPC ensures that the communication cost measured in this thesis reflects the kind
of overhead expected in a realistic microservice inference pipeline rather than a
hand-crafted minimal transport.
The maximum gRPC message size is set to 16~MiB in all stages, which accommodates the
largest intermediate activation tensor in the study (approximately 784~KiB for a split
after \texttt{layer1} in FP32).

The framework supports variable chain depths: the same code path handles two-service
chains used in RQ1.1, RQ1.2, and RQ2, and chains of up to five services used in RQ1.4
and RQ1.5.
The independent variable in each stage is therefore the choice of split point or the
number of chained services, not the communication mechanism.

% ─────────────────────────────────────────────────────────────────────────────
\section{Shared Benchmark Contract}
\label{sec:methods-contract}

To support direct comparison across stages, the thesis defines a shared benchmark
contract that is preserved unless a stage explicitly states otherwise.
The fixed subject under test is ResNet-18 with pretrained weights, executed on CPU in
FP32 with input shape $1 \times 3 \times 224 \times 224$, communicating over gRPC on
port 50051 with a maximum message size of 16~MiB.

The shared timing protocol consists of 50 warmup iterations followed by 200 measured
iterations per condition execution, with a five-second cooldown between successive
conditions, a fixed random seed of 42 for condition ordering within each round, and a
warmup calibration check using a trailing window of 10 iterations and a
coefficient-of-variation (CV) threshold of 0.02.
Thesis-facing stages aggregate five rounds or five paired passes, yielding 1,000 measured
iterations per condition.
The warmup calibration procedure is not used to terminate warmup adaptively: its purpose
is to provide evidence that the fixed 50-iteration warmup is sufficient for stable
measurement in each execution environment.

The shared functional validation contract requires local segment validation before
timing and live gRPC chain validation for any split or Kubernetes deployment, using an
absolute tolerance of $\text{atol} = 1 \times 10^{-5}$ and five validation inputs.
A condition is not benchmarked unless both checks pass.
In practice, all thesis-facing frozen runs achieve a maximum absolute difference of 0.0
against the monolithic reference, confirming that no numerical drift is introduced by
the partition or transport layer.

This shared contract is the invariant across the staged pipeline: what changes between
stages is the infrastructure and security context, not the model, the measurement
discipline, or the validation requirements.

% ─────────────────────────────────────────────────────────────────────────────
\section{CPU Stabilisation}
\label{sec:methods-cpu}

Systematic CPU stabilisation is applied in all local and Kubernetes stages to reduce
variance attributable to hardware noise rather than to the partition boundary.
The following controls are requested in every thesis-facing local run.
Thread counts for PyTorch intra-op and inter-op threads, and for OpenMP, MKL, and
OpenBLAS, are all set to one, ensuring single-threaded model execution.
CPU affinity is pinned to one physical core with simultaneous multithreading (SMT)
excluded, preventing the benchmark process from migrating between cores or sharing a
physical core with a sibling hyperthread.
Process scheduling priority is raised using \texttt{nice~-5} where the runtime permits.
The CPU frequency governor is set to \texttt{performance} and turbo boost is disabled
to prevent frequency scaling from affecting measurements.

These controls are applied symmetrically across all conditions within a benchmark so
that differences in measured latency can be attributed to the independent variable
rather than to hardware variability.

In containerised and cloud stages, some controls are partially unavailable.
In the local \texttt{kind} Kubernetes environment (RQ1.4), governor and turbo-disable
writes to \texttt{/sys} are blocked by the read-only sysfs exposed inside the container;
however, the detected governor is already \texttt{performance} and turbo is already
disabled, so the intended state is achieved through the host configuration.
In AKS (RQ1.5 and RQ2), the \texttt{nice~-5} request is denied due to insufficient
privileges and the benchmark runs at \texttt{nice~=~0}; governor and turbo controls
are unavailable through the managed node interface.
These constraints are documented per stage and treated as realistic limitations of
cloud-native deployment rather than as experimental flaws, because they affect all
conditions within a stage symmetrically.

% ─────────────────────────────────────────────────────────────────────────────
\section{Measurement and Statistical Analysis}
\label{sec:methods-statistics}

\subsection{Primary Metrics}

End-to-end latency is the primary measurement in all stages, defined as the wall-clock
time for one full inference request from the moment the benchmark client dispatches the
input to the moment it receives the final classification output.

Results are summarised using the arithmetic mean, median, and standard deviation of
per-iteration latencies.
The median provides a robust central tendency estimate in the presence of tail spikes.
The mean is used as the primary comparison value because the thesis is concerned with
expected per-request cost.
The 95th-percentile (p95) latency is reported in the deployment and security stages to
characterise tail behaviour under more realistic conditions.

Overhead is reported both in absolute milliseconds and as a percentage of the relevant
baseline mean latency.
In stages that compare conditions across environments (RQ1.5), within-environment
normalised overhead is the primary comparison metric rather than raw absolute latency,
because absolute values are environment-specific and do not transfer across different
host CPUs, kernels, and virtualisation layers.

Several derived metrics are used in specific stages.
The boundary-crossing interval captures the split-benchmark cost associated with
crossing the service boundary, reflecting a composite of serialisation, transport, and
the work executed on the remote side after the boundary.
The activation-transfer burden is derived from the model architecture as the
intermediate tensor size at a split boundary (or cumulative tensor size across all
boundaries in a chain), not measured by a separate network profiler.
The inferred non-compute/platform cost is a derived residual used in the Kubernetes,
AKS, and mTLS analyses to separate platform and communication overhead from model
computation analytically.
In RQ2.1, resource and operational overhead metrics are also reported, including sidecar
CPU and memory, total pod resource deltas, schedule-to-ready time, and the count of
added mesh and security objects.
In RQ2.2, sequential throughput is derived from mean latency as requests per second to
express the throughput impact of confidential-node placement.

\subsection{Cross-Round Consistency}

Cross-round consistency is assessed using the round CV, computed as the standard
deviation of per-round condition means divided by the grand mean.
A low round CV indicates stable execution across independent rounds and supports
treating the reported condition means as representative rather than as artefacts of a
single favourable or unfavourable run.
Cross-round stability and parity validation are treated as the primary indicators of
methodological validity.

\subsection{Effect Size and Significance Tests}

In stages with 1,000 per-condition measured iterations, Cohen's~$d$ is computed from
pooled per-iteration data and the Mann-Whitney $U$ test is reported for completeness.
At $N = 1{,}000$ per condition, $p$-values are near-zero for any non-trivial difference
and should not be over-interpreted as evidence of practical significance.
Cohen's~$d$ provides the more informative measure: values below 0.2 are negligible,
0.2--0.5 are small, 0.5--0.8 are medium, and above 0.8 are large.

% ─────────────────────────────────────────────────────────────────────────────
\section{Carry-Forward and Selection Logic}
\label{sec:methods-carryforward}

The staged pipeline uses a carry-forward mechanism for the local split-selection stages
(RQ1.1 and RQ1.2).
After each of these stages, one or two candidates are selected to enter the next stage
according to two predeclared rules.

The first rule concerns compute degeneracy.
A split condition is classified as compute-degenerate if fewer than 10\% of the total
monolithic compute is allocated to the smaller of the two resulting services.
Compute-degenerate conditions are excluded from main-candidate eligibility because they
produce structurally unbalanced decompositions that do not constitute meaningful
microservice partitions, even if their raw latency is low.

The second rule resolves near-ties among non-degenerate candidates.
When the lowest and second-lowest raw latencies among eligible conditions fall within a
5\% near-best window, the condition with the better structural balance -- lower
activation-transfer burden relative to its compute distribution -- is selected as the
main carry-forward candidate, and the other is retained as the reference candidate.

Carry-forward is not applicable in RQ1.4 and RQ1.5, which benchmark a predefined chain
family rather than selecting among candidates.
It is also not applicable in the RQ2 stages, which fix the deployment topology
established by RQ1 and vary only the security configuration.

% ─────────────────────────────────────────────────────────────────────────────
\section{Stage-Specific Methods}
\label{sec:methods-stages}

\subsection{RQ1.1 -- Coarse Architectural Boundary Selection}
\label{sec:methods-rq11}

RQ1.1 screens the four coarse stage boundaries of ResNet-18 against a monolithic
baseline under the maximally controlled local setup described in
Sections~\ref{sec:methods-cpu} and~\ref{sec:methods-contract}.
The four split conditions place the boundary after \texttt{layer1}, \texttt{layer2},
\texttt{layer3}, and \texttt{layer4}, producing intermediate activation tensors of
approximately 784~KiB, 392~KiB, 196~KiB, and 98~KiB in FP32 respectively.
Both split services run as separate OS processes communicating via gRPC over
\texttt{127.0.0.1:50051}.

The justification for screening at the coarse architectural level before conducting
any finer search is that moving directly into an exhaustive layer-by-layer search would
risk over-fitting the selection to incidental local-execution noise and would not
align with the partitioning literature, which typically evaluates architecturally
meaningful boundaries rather than arbitrary layers~\cite{dnnsplit,deeperthings,survey_partitioning}.

The output of this stage is one main carry-forward candidate and one reference
candidate, selected using the degeneracy and near-best-window rules defined in
Section~\ref{sec:methods-carryforward}.

\subsection{RQ1.2 -- Fine-Grained Refinement}
\label{sec:methods-rq12}

RQ1.2 takes the carry-forward outcome of RQ1.1 as its starting point and introduces one
block-level split point between the two coarse carry-forward anchors.
The new condition, \texttt{split\_after\_layer3\_block0}, bisects the \texttt{layer3}
stage at its internal residual-block boundary.

A key structural property motivates this particular design choice.
\texttt{split\_after\_layer3\_block0} and \texttt{split\_after\_layer3} both produce
intermediate activation tensors of approximately 196~KiB.
The design therefore holds activation-transfer size constant while varying compute
placement, creating a controlled comparison that can isolate the role of compute
distribution independently of communication burden.
This comparison is grounded in the DNN-partitioning literature, which treats compute
distribution as a first-class criterion alongside activation-transfer
cost~\cite{dnnsplit,fine_grained_partitioning}.

The experimental setup and measurement protocol are identical to RQ1.1, ensuring that
the two stages remain directly comparable.

\subsection{RQ1.3 -- Explanatory Synthesis}
\label{sec:methods-rq13}

RQ1.3 is not a split benchmark.
It is a hybrid analytical stage that explains the patterns observed in RQ1.1 and RQ1.2
using two structural properties: activation-transfer size and compute distribution.
The explanatory analysis draws on three evidence sources: the split benchmark summaries
from RQ1.1, the split benchmark summaries from RQ1.2, and a separately conducted
internal timing experiment on the monolithic ResNet-18 model.

The internal timing experiment instruments the monolithic forward pass at the stage,
block, and selected operation level.
It uses the same single-threaded CPU-stabilised environment as the split benchmarks,
with 10 rounds, 50 warmup iterations per round, and 200 measured iterations per round.
Functional equivalence against the uninstrumented model is verified
(\texttt{max\_abs\_diff}~$= 0.0$), and instrumentation overhead is measured explicitly
at approximately 0.645~ms (0.88\% of total forward-pass time), confirming that
profiling does not materially distort the timing distribution.

The justification for this supporting experiment is that split-benchmark
boundary-crossing intervals reflect a composite of serialisation, transport, and
remote-side computation.
Disentangling those components -- particularly distinguishing activation-transfer cost
from compute-placement effects -- requires a direct measurement of local per-block
execution cost that the split benchmark alone cannot provide.

\subsection{RQ1.4 -- Multi-Microservice Kubernetes Chain}
\label{sec:methods-rq14}

RQ1.4 moves from local process-based communication to real Kubernetes in-cluster
services.
The environment is a single-node \texttt{kind} cluster running on Ubuntu/Linux, with
all service pods colocated on the same node.
This colocation rule is applied deliberately so that inter-service gRPC traffic remains
in-cluster rather than traversing an external network, meaning that any observed
overhead reflects orchestration and pod-to-pod communication rather than wide-area
latency.

The stage benchmarks a predefined family of five conditions from
\texttt{monolithic\_k8s\_1svc} to \texttt{chain\_5svc}, corresponding to one through
five synchronously chained ResNet-18 service segments.
The chain family is predefined rather than selected by carry-forward because the purpose
of this stage is not to choose among candidates but to characterise the full cost curve
as a function of service count.
Service pods are allocated one CPU core and 512~MiB of memory with Guaranteed QoS to
prevent throttling or migration during measurement windows.

The use of a single-node \texttt{kind} cluster is an acknowledged limitation: it
eliminates inter-node network hops and does not capture the scheduling and networking
variability of a multi-node or managed cluster.
This limitation is addressed by the transfer-validation design of RQ1.5.

\subsection{RQ1.5 -- AKS Transfer Validation}
\label{sec:methods-rq15}

RQ1.5 re-runs the same predefined chain family from RQ1.4 on a single-node AKS cluster
(\texttt{Standard\_D8s\_v3}, region \texttt{swedencentral}, \texttt{kubenet}
networking).
The stage does not aim to replicate the local Kubernetes absolute latency values, which
are expected to differ because the two environments use different host CPUs, kernels,
virtualisation layers, container runtimes, and networking paths.

The transfer claim is therefore directional rather than numerical: if the same
architectural chain ordering is preserved in AKS, that constitutes evidence that the
ordering reflects the architecture of the model partition rather than incidental
local-cluster noise~\cite{benchmarking_sc15}.
Within-environment normalised overhead is used as the primary comparison metric for
this reason.

All service pods and the benchmark client are colocated on the same AKS node and run
with Guaranteed QoS, applying the same resource and placement discipline as RQ1.4.

\subsection{RQ2.1 -- Service Identity and Mutual TLS Hardening}
\label{sec:methods-rq21}

RQ2.1 introduces a paired execution design.
Each paired pass executes both the plain AKS condition and the mTLS/AuthZ condition in
alternating order, with the order determined by a seeded plan in which three passes run
mTLS first and two passes run plain first.
This alternation reduces the risk that the observed delta is explained by slow
performance drift over the course of the run rather than by the hardening mechanism
itself.

The plain condition uses standard Kubernetes service-to-service communication without
mesh injection.
The mTLS condition adds one Istio sidecar proxy per inference service, service accounts,
workload-level peer-authentication objects, and authorisation-policy objects.
The benchmark client remains outside the mesh; the stage therefore measures hardening
of the inter-service path inside the inference chain, not external ingress hardening.

The managed AKS Istio add-on (\texttt{asm-1-29}) is used as the mesh implementation.
This choice reflects the thesis's focus on realistic cloud-native deployment: managed
add-ons represent the configuration path most likely to be used by practitioners
adopting mTLS in AKS and their behaviour reflects actual cloud-operator defaults rather
than custom tuning.

Security validation is performed for each paired run using four methods: static manifest
validation of peer-authentication and authorisation-policy objects; runtime validation
of policy object counts; a full-chain positive probe confirming that valid upstream-to-
downstream requests succeed; and a direct non-mesh denial probe confirming that
unauthenticated access to protected downstream services is blocked.

The primary topology is \texttt{chain\_2svc} (split after \texttt{layer2}), selected
because it is the lowest-overhead non-monolithic chain in the preceding stages.
A \texttt{chain\_5svc} stress topology is included to determine how mTLS overhead
grows when more service boundaries are protected.

A known limitation of sidecar-based mTLS is that it protects the inter-proxy
communication path but does not prevent plaintext from existing inside the pod boundary
after TLS termination.
This limitation is acknowledged explicitly: RQ2.1 is claimed as a communication-level
hardening layer, not as a complete inter-service secrecy mechanism~\cite{poudel2025mazu}.

\subsection{RQ2.2 -- Confidential VM Execution Hardening}
\label{sec:methods-rq22}

RQ2.2 keeps the communication-security contract from RQ2.1 fixed and changes the
execution environment of the inference services.
The confidential condition places selected services on AMD SEV-SNP confidential node
pools (\texttt{Standard\_DC8as\_v5}); the standard condition uses matched
non-confidential AMD nodes (\texttt{Standard\_D8as\_v5}) in the same AKS cluster.
Fresh paired standard baselines are collected for each RQ2.2 artifact set rather than
reusing older RQ2.1 data, ensuring that incremental overhead can be attributed
unambiguously to the confidential-node placement rather than to session-level drift.

Two configurations are evaluated.
The selective TEE run places only the downstream service (\texttt{service2}) on
confidential nodes, isolating the cost of protecting the service that receives
intermediate activations.
The full TEE run places both services on confidential nodes, measuring the cost of
extending the confidential-execution boundary to the entire two-service chain.

The protection boundary of AMD SEV-SNP must be stated precisely.
SEV-SNP provides VM-level memory encryption and integrity protection against an
untrusted host platform, but the guest operating system, application process, PyTorch
runtime, model weights, and Istio sidecar all remain inside the same VM trust boundary.
RQ2.2 therefore uses the term \emph{VM-level confidential execution} throughout rather
than claiming application-level enclave isolation.

\subsection{RQ2.3 -- Security-Hardening Trade-off Synthesis}
\label{sec:methods-rq23}

RQ2.3 is a synthesis stage: no new benchmark data is collected.
The stage combines the plain AKS baseline from RQ1.5, the mTLS/AuthZ results from
RQ2.1, and the selective and full TEE results from RQ2.2 to compare the four hardening
configurations across measured cost, operational complexity, and protection scope.

The synthesis is structured around the principle that the preferred hardening
configuration is threat-model-dependent rather than universally optimal.
Applying every available security mechanism to every service boundary is not always the
best engineering choice, because each additional protection layer adds measurable
latency, resource, and operational cost.
RQ2.3 therefore frames its output as a structured trade-off comparison rather than a
single global ranking, allowing a practitioner to match a security configuration to a
stated threat model and performance budget.

% ─────────────────────────────────────────────────────────────────────────────
\section{Artifact Freezing and Reproducibility}
\label{sec:methods-artifacts}

Thesis-facing runs are exported and frozen as artifact sets before analysis is
performed.
This means that the reported results are derived from a fixed, immutable snapshot of the
raw measurement data rather than from ad hoc terminal output or re-runnable live
scripts.
Each frozen artifact set records the full execution environment: operating system,
Python and PyTorch version, Kubernetes version and cluster name, node type, namespace,
container image digest, condition ordering, and parity validation outcomes.

The separation between the frozen artifact and the analysis code ensures that the
reported numbers remain stable even if benchmark scripts are modified for subsequent
exploratory runs.
The repository contains both thesis-facing configs and older exploratory or smoke-test
configs; the artifact freeze identifies which run is the thesis-facing primary result.

% ─────────────────────────────────────────────────────────────────────────────
\section{Limitations and Threats to Validity}
\label{sec:methods-limitations}

Several methodological limitations should be acknowledged.

\textbf{Single-node clusters.}
Both the RQ1.4 \texttt{kind} cluster and the RQ1.5 AKS cluster are single-node.
Single-node deployments eliminate inter-node network hops and do not capture the
scheduling and networking variability of a multi-node production cluster.
Results therefore reflect in-cluster pod-to-pod latency on a single physical or virtual
machine and are not directly generalisable to distributed multi-node deployments.

\textbf{CPU-only inference.}
All stages use CPU inference with a single thread.
GPU-accelerated inference would substantially shift the compute distribution, reduce
per-inference time, and change the relative weight of the gRPC boundary cost.
The results are not directly generalisable to GPU-hosted microservice deployments.

\textbf{Single model and workload.}
The thesis evaluates one model (ResNet-18) at batch size one with one input shape.
Different architectures, larger activation maps, or higher-batch workloads may produce
different cost curves and different optimal split regions.

\textbf{Controlled conditions versus production variability.}
The local stages apply strict CPU stabilisation and achieve low cross-round variance.
The Kubernetes and AKS stages are less controllable and exhibit higher variability.
Neither environment captures the full noise of a shared multi-tenant cloud cluster,
where CPU steal, network congestion, and noisy neighbours can substantially affect tail
latency.

\textbf{mTLS threat-model scope.}
Sidecar-based mTLS protects the inter-proxy communication path but does not prevent
plaintext from existing inside the endpoint after TLS termination.
The security claims for RQ2.1 are explicitly scoped to communication-level hardening.

\textbf{SEV-SNP threat-model scope.}
AMD SEV-SNP protects the VM boundary against the host platform but not against a
compromised guest OS or application process, and several side-channel and availability
threats fall outside its core guarantee.
RQ2.2 claims VM-level confidential execution, not complete application-level isolation.

% ─────────────────────────────────────────────────────────────────────────────
\section{Ethical Considerations}
\label{sec:methods-ethics}

This thesis does not involve human participants, personal data, or biological material.
The study is a software systems benchmark using a publicly available pretrained image
classification model and standard synthetic inputs; no classification decisions derived
from this model are used to affect any real-world outcomes, and no personal images or
sensitive datasets are processed.

ResNet-18 with pretrained ImageNet weights inherits any biases or limitations of the
ImageNet training process.
However, since the thesis does not evaluate model accuracy and does not deploy the model
in a decision-making context, these concerns are not directly relevant to the research
questions addressed here.

Cloud compute resources are used responsibly: experiments are scoped to the minimum
infrastructure required to answer each research question, and resources are
deallocated after each benchmark stage.
All measurement artifacts contain only timing data and system metadata; no sensitive
data is stored on cloud infrastructure.
